\subsection{Algebraically disjoint extensions}

DEFINITION 5. --- Let L be an extension of a field K, and E and F two subextensions of L. Then E and F are said to be algebraically disjoint (over K) and E is said to be algebraically disjoint from F over K, if for every subset A (resp. B) of E (resp. F) algebraically free over K, A and B are disjoint and A ∪ B is algebraically free over K.

Remarks. --- 1) If E is a subextension of L which is algebraic over K, then it is algebraically disjoint from every subextension F of L. For an extension of K to be algebraic it is necessary and sufficient that it should be algebraically disjoint from itself.

\begin{enumerate}
\def\labelenumi{\arabic{enumi})}
\setcounter{enumi}{1}
\item
  It may happen that E is algebraically disjoint from F over K, but not over a subfield \(K_{0}\) of K. * For example C is algebraically disjoint from itself over R but not over Q. *
\item
  It is clear that if E is algebraically disjoint from F over K, when E and F are considered as subextensions of L, the same is true when they are considered as subextensions of \(\mathbf{K}(\mathbf{E} \cup \mathbf{F})\) and conversely.
\end{enumerate}

PROPOSITION 11. --- If E and F are algebraically disjoint over K then E ∩ F is algebraic over K.

This follows from Def. 5.

PROPOSITION 12. --- Let L be an extension of a field K and E, F subextensions of L. Then the following conditions are equivalent :

\begin{enumerate}
\def\labelenumi{\alph{enumi})}
\item
  E and F are algebraically disjoint;
\item
  there exists a transcendence basis of E over K which is algebraically free over F;
\item
  every subset of E which is algebraically free over K is algebraically free over F. Let us introduce the following conditions:
\end{enumerate}

\(b^{\prime})\) there exists a transcendence basis of \(\mathbf{F}\) over \(\mathbf{K}\) which is algebraically free over \(\mathbf{E}\) ;

\(c^{\prime})\) every subset of \(\mathbf{F}\) which is algebraically free over \(\mathbf{K}\) is algebraically free over \(\mathbf{E}\) .

\(a) \Rightarrow b')\) : Suppose that E and F are algebraically disjoint. Let B (resp. C) be a transcendence basis of E (resp. F) over K. Then \(B \cap C = \emptyset\) and \(B \cup C\) is algebraically free over K, hence C is algebraically free over \(K(B)\) (V, p.~107, Prop. 4); since E is algebraic over \(K(B)\) , Prop. 6 of V, p.~108 shows C to be algebraically free over E.

\(b') \Rightarrow c)\) : Suppose that there exists a transcendence basis C of F over K which is algebraically free over E. Let A be a subset of E algebraically free over K. Then C is algebraically free over K(A), hence A is algebraically free over K(C) (V, p.~107, Prop. 4) and therefore over F (V, p.~108, Prop. 6), because F is algebraic over K(C).

\(c) \Rightarrow a)\) : This follows at once from Prop. 4 (V, p.~107).

The implications \(a) \Rightarrow b) \Rightarrow c') \Rightarrow a)\) may be proved in the same way.

COROLLARY. --- Let E and F be algebraically disjoint over K. Let E' be the relative algebraic closure of E in L and F' that of F (V, p.~19). Then E' and F' are algebraically disjoint over K.

Let B be a transcendence basis of E over K; this is also a transcendence basis of \(E'\) over K. Since E is algebraically disjoint from F over K, B is algebraically free over F, hence over \(F'\) (V, p.~108, Prop. 6); now we can apply Prop. 12.

PROPOSITION 13. --- Let L be an extension of a field K and E, F two subextensions of L.

\begin{enumerate}
\def\labelenumi{\alph{enumi})}
\tightlist
\item
  We have \(\operatorname{tr} \cdot \deg_{\mathbb{F}} \mathrm{F}(\mathrm{E}) \leqslant \operatorname{tr} \cdot \deg_{\mathbb{K}} \mathrm{E}\) . When \(\mathrm{E}\) and \(\mathrm{F}\) are algebraically disjoint over \(\mathbf{K}\) , then every transcendence basis of \(\mathrm{E}\) over \(\mathbf{K}\) is a transcendence basis of
\end{enumerate}

\(F(E)\) over F and we have tr. \(\deg_{F}F(E) = \text{tr. } \deg_{K}E\) . Conversely, this equality implies that E and F are algebraically disjoint over K when tr. \(\deg_{K}E\) is finite.

\begin{enumerate}
\def\labelenumi{\alph{enumi})}
\setcounter{enumi}{1}
\item
  We have \(\operatorname{tr} \cdot \deg_{\mathbf{K}} \mathrm{K}(\mathrm{E} \cup \mathrm{F}) \leqslant \operatorname{tr} \cdot \deg_{\mathbf{K}} \mathrm{E} + \operatorname{tr} \cdot \deg_{\mathbf{K}} \mathrm{F}\) . When \(\mathrm{E}\) and \(\mathrm{F}\) are algebraically disjoint over \(\mathbf{K}\) , we have \(\operatorname{tr} \cdot \deg_{\mathbf{K}} \mathrm{K}(\mathrm{E} \cup \mathrm{F}) = \operatorname{tr} \cdot \deg_{\mathbf{K}} \mathrm{E} + \operatorname{tr} \cdot \deg_{\mathbf{K}} \mathrm{F}\) . Conversely, this equality implies that \(\mathrm{E}\) and \(\mathrm{F}\) are algebraically disjoint over \(\mathbf{K}\) when \(\mathrm{E}\) and \(\mathrm{F}\) are of finite transcendence degree over \(\mathbf{K}\) .
\item
  Let B be a transcendence basis of E over K; then E is algebraic over K(B), and Cor. 2 of V, p.~18 shows that F(E) is algebraic over F(K(B)) = F(B). By Th. 2 (V, p.~109) B contains a transcendence basis of F(E) over F; when E is algebraically disjoint from F over K, B is algebraically free over F (Prop. 12) and this is then a transcendence basis of F(E) over F. The three first assertions of a) follow from this. Suppose now that E is of finite transcendence degree over K, equal to that of F(E) over F; since F(E) is algebraic over F(B) and Card B = tr . deg\_F F(E), Cor. 1 of V, p.~110 shows B to be algebraically free over F, so E is algebraically disjoint from F over K (Prop. 12).
\item
  We have \(\mathbf{K}(\mathbf{E} \cup \mathbf{F}) = \mathbf{F}(\mathbf{E})\) and so the Cor. of V, p.~110 implies the equality:
\end{enumerate}

\[
\operatorname{tr} \cdot \deg_ {K} K (E \cup F) = \operatorname{tr} \cdot \deg_ {F} F (E) + \operatorname{tr} \cdot \deg_ {K} F.
\]

Now \(b)\) follows at once from \(a)\) and this equality.

PROPOSITION 14. --- Let L be an extension of a field K, E and F two subextensions of L and B a transcendence basis of E over K. For E and F to be algebraically disjoint over K it is necessary and sufficient that K(B) and F should be linearly disjoint over K.

For E and F to be algebraically disjoint over K it is necessary and sufficient for B to be algebraically free over F (Prop. 12), that is, for the monomials in the elements of B to be linearly independent over F. Since these monomials form a basis of the vector K-space K{[}B{]}, it comes to the same to say that K{[}B{]} and F are linearly disjoint over K. Finally, since K(B) is the field of fractions of K{[}B{]}, Prop. 6 of V, p.~14 shows that K{[}B{]} and F are linearly disjoint if and only if this is so for K(B) and F.

COROLLARY 1. --- If E and F are linearly disjoint, then E is algebraically disjoint from F over K. Conversely, if E is a pure extension of K and is algebraically disjoint from F over K, then E and F are linearly disjoint over K.

COROLLARY 2. --- Every pure extension of K is linearly disjoint from every algebraic extension of K; in particular, K is relatively algebraically closed in every pure extension of K.

